\section{Principi e linee guida}

Nessun progetto è al 100\% originale, anzi molto spesso si trovano delle similitudini con altri progetti che permettono di definire delle indicazioni valide, più o meno generali, più o meno obbligatorie, che possono aiutare il progettista.

\begin{tabularx}{\textwidth}{>{\centering\arraybackslash}X|>{\centering\arraybackslash}X|>{\centering\arraybackslash}X|}
                           & \bld{Bassa coercitività} & \bld{Alta coercitività} \\
    \hline
    \bld{Bassa generalità} & Linee guida              & Regole di progetto      \\
    \hline
    \bld{Alta generalità}  & Principi                 & Standard
\end{tabularx}

\subsection{Standard della human-system interaction}

Preparati dalla ISO, gli standard dell'ingegneria dell'usabilità fanno capo al \itc{TC 159 - Ergonomics}, diviso in quattro \itc{Sub-Committee}:

\begin{itemize}
    \item TC 159/SC 1 - General ergonomics principles
    \item TC 159/SC 3 - Anthropometry and biomechanics
    \item TC 159/SC 4 - Ergonomics of human-system interaction
    \item TC 159/SC 5 - Ergonomics of the physical environment
\end{itemize}

\subsubsection{TC 159/SC 4: Standardizzazione ergonomica}

La sua missione è quella di standardizzare l'ergonomia dell'interazione fra sistemi e le persone che li progettano, fabbricano, usano e mantengono.

Produce i seguenti standard:

\begin{itemize}
    \item \itc{ISO-13407 Human-centred design processes for interactive systems}: ha lo scopo di aiutare coloro che hanno l responsabilità di gestire i processi di progettazione di hardware e software a pianificare in modo adeguato le attività di progettazione human-centred.
    \item \itc{ISO-9241}: Si tratta dello standard principale relativo alla human-computer interaction. È molto ampio, ed è composto da numerosi documenti separati, in evoluzione da una ventina d’anni.
    \item \itc{ISO-14915 Software ergonomics for multimedia user-interfaces}, composto da tre documenti:
          \begin{enumerate}
              \item Design principles and framework
              \item Multimedia navigation and control
              \item Media selection and combination
          \end{enumerate}
\end{itemize}

\subsection{I principi del dialogo}

Espressi nella \itc{ISO-9241}, sono i sette principi che ogni dialogo utente-macchina deve possedere. Non sono indipendenti, e a volte bisogna trovare un compromesso tra questi sette. I principi possono essere usati come una misura della qualità e dell'usabilità di un sistema.

\subsubsection{Adeguatezza al compito}

Afferma che le funzionalità e le interazioni del sistema devono essere modellate sulle caratteristiche del compito che l'utente deve svolgere, e non sulla tecnologia e sulle caratteristiche del sistema in sé. Abbiamo sette specifiche di adeguatezza:

\bld{Dialogo adeguato al compito}: i passi del dialogo devono essere adeguati al compito da svolgere, il dialogo deve assegnare al sistema tutto ciò che è automatizzabile, e le operazioni assegnate all'utente devono essere eseguite con facilità.

\bld{Informazione adeguata al compito}: il sistema deve dare all'utente le informazioni utili per svolgere un compito, tenendo a mente la capacità della memoria a breve termine (vedi \autoref{chap:memoria_breve_termine}).

\bld{Dialogo essenziale}: il sistema non deve presentare informazioni ridondanti all'utente.

\bld{Dispositivi d'input e output adeguati al compito}: scegliere sempre il miglior tipo di dispositivo per l'input e l'output (anche più di uno in caso di multi-modalità).

\bld{Formati d'input e output adeguati al compito}: i formati dei dati devono essere adeguati al compito.

\bld{Default tipici}: scegliere i valori di default che riflettano le scelte più comuni degli utenti.

\bld{Compatibilità con i documenti}: disporre le informazioni video in input e in output quanto più vicine alla loro rappresentazione cartacea.

\subsubsection{Auto-descrizione}

Il sistema deve comunicare all'utente in ogni momento cosa può fare, come, e cosa sta accadendo.
